\documentclass[11pt,a4paper]{ctexart}
\usepackage[margin=1.55cm,top=1.6cm,bottom=1.6cm]{geometry}
\usepackage{enumitem}
\usepackage{titlesec}
\usepackage{xcolor}
\usepackage{hyperref}
\usepackage{tabularx}
\usepackage{fontawesome5}

\setmainfont{Times New Roman}
\definecolor{accent}{HTML}{2D6EA7}
\definecolor{textgray}{HTML}{5B6875}
\definecolor{rulegray}{HTML}{D2DCE6}
\hypersetup{colorlinks=true,urlcolor=accent,linkcolor=accent}
\setlength{\parindent}{0pt}
\setlength{\parskip}{0pt}
\linespread{1.03}
\setlist[itemize]{leftmargin=1.3em,label=\textcolor{accent}{\footnotesize\textbullet},itemsep=3.6pt,topsep=4.2pt,parsep=0pt}
\pagestyle{empty}

\titleformat{\section}
  {\bfseries\fontsize{11pt}{13pt}\selectfont\color{accent}}
  {}{0pt}{}
  [\vspace{0.35em}\color{rulegray}\titlerule]
\titlespacing*{\section}{0pt}{1.4em}{0.6em}

\newcommand{\sectionicon}[1]{\raisebox{-0.08em}{\small #1}\hspace{0.45em}}

\newcommand{\educationheader}[2]{%
  \begin{tabularx}{\textwidth}{@{}Xr@{}}
    {\bfseries #1} & {\small\color{textgray}#2}\\[-0.12em]
  \end{tabularx}%
}

\newcommand{\educationline}[1]{%
  \noindent{\small #1}\par\vspace{3pt}%
}

\newcommand{\research}[3]{%
  \begin{tabularx}{\textwidth}{@{}Xr@{}}
    {\bfseries #1} & {\small\color{textgray}#2}
  \end{tabularx}%
}

\newcommand{\dateditem}[2]{%
  \begin{tabularx}{\textwidth}{@{}Xr@{}}
    {#1} & {\small\color{textgray}#2}\\[-0.12em]
  \end{tabularx}\vspace{0.34em}%
}

\begin{document}

% ===== Header =====
\begin{center}
  {\bfseries\fontsize{22pt}{26pt}\selectfont 王鹤宇}\\[0.48em]
  \small\color{textgray}
  \href{mailto:wangheyu23@mails.ucas.ac.cn}{\faEnvelope\ wangheyu23@mails.ucas.ac.cn}
  \hspace{1.8em}
  \href{https://github.com/heyu-wang}{\faGithub\ github.com/heyu-wang}
\end{center}
\vspace{0.1em}

% ===== Education =====
\section[教育背景]{\sectionicon{\faGraduationCap}教育背景}
\educationheader{中国科学院大学}{2023 -- 至今}
\vspace{0.32em}
\educationline{计算机科学与技术专业本科生}
\educationline{\textbf{GPA：}3.85}
\educationline{\textbf{专业排名：}17/86}
\educationline{\textbf{CET-6：}558}
{\small \textbf{课程：}自然语言处理（93）；人工智能计算系统（96）；操作系统实验（93）}

% ===== Skills =====
\section[技术能力]{\sectionicon{\faTools}技术能力}
{\small
\textbf{编程语言：}C、Assembly、Python、Verilog。\\[3.8pt]
\textbf{工具：}SSH、Git、Docker、\LaTeX。
}

% ===== Research Experience =====
\section[研究经历]{\sectionicon{\faFlask}研究经历}

\research{操作系统实验}{2025 年秋季}{https://github.com/heyu-wang/OS}
\begin{itemize}
  \item 基于 RISC-V 开发操作系统内核，完成 6 个实验任务及进阶扩展。
  \item 实现进程调度、虚拟内存、文件系统和多核 SMP 等核心服务。
  \item 完成中断处理、系统调用及 e1000 网卡驱动，实现网络通信。
\end{itemize}

\vspace{0.5em}
\research{计算机体系结构实验与计算机组成原理实验}{2025 年秋季}{https://github.com/heyu-wang/CA}
\begin{itemize}
  \item 在计算机组成原理实验中实现基于 MIPS 和 RISC-V 的多周期 CPU。
  \item 在体系结构实验中协作开发基于 LoongArch 的五级流水线 CPU。
  \item 实现基本框架、Cache 和 AXI 总线，并完成与流水线的集成。
\end{itemize}

\vspace{0.5em}
\research{大语言模型越狱自动研究与评测}{2026 年夏季 -- 至今}{}
{\small\itshape\color{textgray}合作项目}\par
\begin{itemize}
  \item 参与分析 JailbreakBench、StrongREJECT 和 HarmBench 的 LLM-as-a-Judge 评测协议。
  \item 在基于 Harbor 的合作研究框架中，参与 JailbreakBench 接入及黑盒 autoresearch 任务配置。
  \item 支持受查询预算约束、隐藏测试集隔离验证和开发集结果记录的可复现实验流程。
\end{itemize}

% ===== Awards =====
\section[荣誉奖项]{\sectionicon{\faAward}荣誉奖项}
\dateditem{校级三好学生}{2024、2025、2026}
\dateditem{校级二等奖学金}{2023 -- 2024}
\dateditem{校级三等奖学金}{2024 -- 2025}

% ===== Public Affairs =====
\section[学生工作]{\sectionicon{\faUsers}学生工作}
\dateditem{中国科学院大学学生会生活文化部干事}{2023 -- 2024}
\dateditem{中国科学院大学星空志愿者协会外联部部长}{2024 -- 2025}

\end{document}
